% ----------------------------------------------------------
% Justificativa
% ----------------------------------------------------------
\chapter{JUSTIFICATIVA} \label{cha:fundamentacao}

Esse cenário de crescimento global, descrito anteriormente, também se reflete no mercado nacional, que se consolidou como um ecossistema dinâmico de entretenimento e se tornou uma fonte relevante de receita e geração de empregos qualificados na economia brasileira \cite{neto2023}. Sendo assim, promover a interação entre desenvolvedores nacionais e sua comunidade de jogadores é algo que pode movimentar ainda mais esse mercado e impulsionar a inovação no setor.

Entretanto, a dispersão do feedback entre diferentes ambientes representa um desafio significativo também para os desenvolvedores brasileiros. Espaços genéricos, como o Reddit, são amplamente utilizados para o compartilhamento de opiniões sobre jogos, mas sua natureza heterogênea faz com que ideias relevantes se percam em meio a conteúdos diversos, como memes e outras publicações sem relação direta com o desenvolvimento do produto. Essa fragmentação torna mais custoso para o desenvolvedor reunir e interpretar as contribuições da comunidade de forma contínua ao longo do projeto, dificultando o aproveitamento do feedback como fonte de informação para suas decisões criativas e técnicas.

Diante dessa problemática, a implementação de um ambiente especializado para centralizar ideias e feedback associados a jogos mostra-se fundamental, não apenas por impulsionar a inovação no setor, mas também por estimular a cocriação entre os membros da comunidade. Sua relevância reside na otimização do processo de busca por insights, eliminando a necessidade de filtrar conteúdos irrelevantes que atualmente se encontram misturados a publicações sem relação com o tema.

Nesse sentido, a proposta do protótipo de interface para essa plataforma supera a simples organização de dados, posicionando-se como a etapa inicial de um mecanismo estratégico para otimizar o ciclo de desenvolvimento de jogos digitais no ecossistema nacional. Ao estruturar visualmente um ambiente voltado ao crowdsourcing, este trabalho busca demonstrar a viabilidade de uma interface capaz de auxiliar futuramente os desenvolvedores na validação de requisitos e na identificação de funcionalidades prioritárias, contribuindo para a redução de incertezas durante o processo de desenvolvimento.
